\chapter{Cómo se forma la imagen: píxeles y geometría de barrido}\label{cap:escaneo}

\objetivos{
\begin{itemize}
  \item Describir cómo un sensor de barrido transversal construye una imagen línea a línea.
  \item Explicar por qué los píxeles crecen hacia los bordes de la franja y calcular cuánto.
  \item Explicar el efecto pajarita y por qué produce detecciones repetidas del mismo fuego.
  \item Describir cómo VIIRS limita el crecimiento de sus píxeles y qué consecuencias tiene.
  \item Interpretar las columnas \texttt{scan} y \texttt{track} de los datos de focos.
\end{itemize}}

\section{Un espejo que gira}

MODIS y VIIRS no tienen una cámara con millones de sensores como la de un teléfono. Usan unos
pocos detectores y un espejo o telescopio giratorio. En cada giro, el instrumento recorre con la
mirada una línea del suelo perpendicular a la trayectoria del satélite, de un borde de la franja
al otro, tomando muestras a intervalos regulares. Mientras tanto el satélite avanza, de modo que
el giro siguiente recorre una línea un poco más adelante. La imagen se va tejiendo así, línea a
línea. Este diseño se llama \termino{barrido transversal} o \emph{whiskbroom}
(figura~\ref{fig:whisk}).

\begin{figure}[ht]
\centering
\begin{tikzpicture}[scale=0.92]
  \fill[bosque!18] (-6.2,-0.6)--(6.2,-0.6)--(7.6,1.6)--(-4.8,1.6)--cycle;
  \foreach \i in {0,1,2,3}{
    \fill[brasa!\the\numexpr25+15*\i\relax] (-5.9+0.2*\i,-0.4+0.4*\i)--(6.0+0.2*\i,-0.4+0.4*\i)--(6.2+0.2*\i,-0.05+0.4*\i)--(-5.7+0.2*\i,-0.05+0.4*\i)--cycle;
  }
  \begin{scope}[shift={(0.5,6)}]
    \fill[gray!45] (-0.45,-0.28) rectangle (0.45,0.28);
    \fill[cielo!70] (-1.7,-0.22) rectangle (-0.55,0.22); \fill[cielo!70] (0.55,-0.22) rectangle (1.7,0.22);
  \end{scope}
  \draw[->,very thick,ceniza] (2.6,6.0)--(3.9,6.6) node[right,font=\footnotesize,align=left]{avance del\\ satélite};
  \draw[fuego,thick] (0.5,5.7)--(-5.2,1.0); \draw[fuego,thick] (0.5,5.7)--(6.7,1.0);
  \draw[fuego,thick,dashed] (0.5,5.7)--(0.7,1.0);
  \draw[fuego,very thick,->] (-2.2,3.3) arc (215:325:3.3 and 1.1);
  \node[font=\footnotesize,fuego] at (1.6,2.25) {barrido};
  \node[font=\footnotesize] at (0.7,0.55) {nadir};
  \node[font=\footnotesize,align=center] at (-5.6,2.2) {borde de\\ la franja};
  \draw (0.5,4.6) arc (270:309:1.1); \node[font=\footnotesize] at (1.35,4.3) {$\theta$};
  \draw[<->] (-5.7,-0.9)--(6.2,-0.9) node[midway,below,font=\footnotesize]{ancho de la franja};
\end{tikzpicture}
\caption{Barrido transversal. El espejo recorre una línea de borde a borde; el avance del
satélite desplaza la línea siguiente. El ángulo $\theta$ entre la vertical y la dirección de
mirada es el \emph{ángulo de barrido}.}
\label{fig:whisk}
\end{figure}

Hay dos direcciones que conviene nombrar desde ya:

\begin{itemize}
  \item \textbf{a lo largo del barrido} (\emph{along-scan}): perpendicular al avance del
  satélite, la dirección en la que gira el espejo;
  \item \textbf{a lo largo de la traza} (\emph{along-track}): la dirección de avance.
\end{itemize}

\section{Qué es exactamente un píxel}

Cada detector ve, en cada instante, un cono muy estrecho de direcciones.

\begin{definicion}{campo de visión instantáneo (IFOV)}
El campo de visión instantáneo es el ángulo, medido desde el sensor, que abarca un detector en
un instante. Es una propiedad fija del instrumento. La porción de suelo que cae dentro de ese
ángulo es la \emph{huella} del píxel.
\end{definicion}

El ángulo es fijo; la huella no. En el nadir, a una altura $h$, un detector con IFOV $\alpha$
(en radianes) abarca en el suelo
\begin{equation}
  p_0 = h\,\alpha .
\end{equation}
Para MODIS, la huella de 1~km a 705~km de altura corresponde a
$\alpha = 1/705 = \SI{1.42e-3}{\radian}$, unos 0,08 grados.

\begin{atencion}[El píxel no es un cuadrado con bordes nítidos]
La respuesta del detector no es uniforme dentro de la huella ni cae a cero justo en su borde.
Además, el espejo se mueve mientras se toma la muestra. En MODIS esto hace que la huella efectiva
a lo largo del barrido sea más ancha que el espaciado entre muestras: las huellas de píxeles
vecinos se \textbf{solapan} en parte. Un fuego pequeño e intenso situado cerca del límite entre
dos píxeles puede elevar la señal de ambos y ser detectado dos veces en la misma línea.
\end{atencion}

\section{Los píxeles crecen hacia los bordes}

Cuando el espejo apunta hacia un lado, el mismo cono angular alcanza el suelo más lejos y,
además, en oblicuo. Ambas cosas agrandan la huella (figura~\ref{fig:crecimiento}).

\begin{figure}[ht]
\centering
\begin{tikzpicture}[scale=1.0]
  \draw[bosque,very thick] (-1,0)--(11.5,0);
  \fill (1,5) circle (0.09) node[above,font=\footnotesize]{sensor};
  \draw[dashed,ceniza] (1,5)--(1,0);
  \fill[fuego,opacity=0.3] (1,5)--(0.72,0)--(1.28,0)--cycle;
  \draw[fuego] (1,5)--(0.72,0) (1,5)--(1.28,0);
  \fill[fuego,opacity=0.3] (1,5)--(5.45,0)--(6.6,0)--cycle;
  \draw[fuego] (1,5)--(5.45,0) (1,5)--(6.6,0);
  \fill[fuego,opacity=0.3] (1,5)--(8.4,0)--(11.2,0)--cycle;
  \draw[fuego] (1,5)--(8.4,0) (1,5)--(11.2,0);
  \draw[very thick,black] (0.72,-0.12)--(1.28,-0.12);
  \draw[very thick,black] (5.45,-0.12)--(6.6,-0.12);
  \draw[very thick,black] (8.4,-0.12)--(11.2,-0.12);
  \node[below,font=\footnotesize] at (1,-0.2) {nadir};
  \node[below,font=\footnotesize] at (6,-0.2) {$\theta$ medio};
  \node[below,font=\footnotesize] at (9.8,-0.2) {$\theta$ grande};
  \draw (1,3.9) arc (270:316:1.1); \node[font=\footnotesize] at (1.65,3.75) {$\theta$};
  \draw[<->] (0.6,5)--(0.6,0) node[midway,left,font=\footnotesize]{$h$};
\end{tikzpicture}
\caption{El mismo ángulo de visión cubre cada vez más suelo conforme la mirada se inclina (vista
de perfil, Tierra plana para simplificar).}
\label{fig:crecimiento}
\end{figure}

Con una Tierra plana, la geometría es elemental. La distancia del sensor al suelo en la
dirección $\theta$ es $h/\cos\theta$, así que la huella crece en ese factor en la dirección de
avance. A lo largo del barrido hay un segundo factor $1/\cos\theta$, porque el cono corta el
suelo en oblicuo:
\begin{equation}\label{eq:plano}
  p_{\text{traza}} \approx \frac{p_0}{\cos\theta},
  \qquad
  p_{\text{barrido}} \approx \frac{p_0}{\cos^2\theta}
  \qquad\text{(Tierra plana).}
\end{equation}

La Tierra, sin embargo, es curva, y cerca de los bordes de una franja de más de 2000~km la
curvatura importa mucho: el suelo <<se aleja>> del sensor más de lo que predice el modelo plano
y la mirada llega aún más rasante. El cálculo exacto usa el triángulo formado por el centro de la
Tierra, el satélite y el punto observado (figura~\ref{fig:curva}). Si $R$ es el radio terrestre,
el ángulo $\beta$ con que la línea de mirada llega al suelo, medido desde la vertical local, se
obtiene con el teorema del seno:
\begin{equation}\label{eq:cenit}
  \sin\beta = \frac{R+h}{R}\,\sin\theta ,
\end{equation}
y las dimensiones de la huella son
\begin{equation}\label{eq:curvo}
  p_{\text{traza}} = p_0\,\frac{d}{h},
  \qquad
  p_{\text{barrido}} = p_0\,\frac{d}{h}\,\frac{1}{\cos\beta},
\end{equation}
donde $d$ es la distancia real del sensor al punto observado.

\begin{figure}[ht]
\centering
% O es el centro común de ambos radios; h y los ángulos se dibujan aparte.
\begin{tikzpicture}[scale=.95,font=\footnotesize]
  \coordinate (O) at (0,0);
  \coordinate (S) at (0,6.3);
  \coordinate (N) at (0,5);
  \coordinate (P) at (67:5);
  \draw[bosque,very thick] (50:5) arc (50:130:5);
  \draw[ceniza,dashed] (O)--(N) node[midway,left] {$R$};
  \draw[ceniza,dashed] (N)--(S);
  \draw[ceniza,dashed] (O)--(P) node[pos=.45,right] {$R$};
  \fill (O) circle (.06) node[below=3pt] {centro de la Tierra};
  \fill (S) circle (.08) node[above=3pt] {satélite};
  \fill (P) circle (.07);
  \node[anchor=west] at ($(P)+(.2,-.15)$) {punto observado};
  \draw[fuego,thick] (S)--(P) node[midway,above right] {$d$};
  \draw[ceniza,dashed] (P)--($(P)+(67:1.5)$);
  \node[anchor=west] at ($(P)+(67:1.65)$) {vertical local};
  \draw[<->] (-.5,5)--(-.5,6.3) node[midway,left] {$h$};
  \draw (0,5.65) arc[start angle=270,end angle=319,radius=.65];
  \node at (.33,5.22) {$\theta$};
  \draw ($(P)+(67:.55)$) arc[start angle=67,end angle=139,radius=.55];
  \node at ($(P)+(103:.85)$) {$\beta$};
\end{tikzpicture}

\caption{Con la Tierra curva, la mirada llega al suelo con un ángulo $\beta$ respecto de la
vertical local que es mayor que el ángulo de barrido $\theta$. Los dos radios parten
del mismo centro terrestre. Esquema geométrico, no a escala.}
\label{fig:curva}
\end{figure}

\begin{ejemplo}{el píxel de MODIS en el borde de la franja}
MODIS barre hasta $\theta = 55^\circ$. Calcula el tamaño de su píxel de 1~km en el borde, con el
modelo plano y con el modelo curvo ($R = \SI{6371}{\kilo\metre}$, $h = \SI{705}{\kilo\metre}$).
\solucion \emph{Modelo plano.} $\cos 55^\circ = 0{,}574$. Entonces
$p_{\text{traza}} = 1/0{,}574 = \SI{1.74}{\kilo\metre}$ y
$p_{\text{barrido}} = 1/0{,}574^2 = \SI{3.04}{\kilo\metre}$.

\emph{Modelo curvo.} De~\eqref{eq:cenit}: $\sin\beta = (7076/6371)\times\sin 55^\circ =
0{,}910$, luego $\beta = 65{,}5^\circ$. El ángulo en el centro de la Tierra es $\beta - \theta =
10{,}5^\circ$, y por el teorema del seno $d = R\,\sin(10{,}5^\circ)/\sin(55^\circ) =
\SI{1414}{\kilo\metre}$. Con~\eqref{eq:curvo}:
\[
 p_{\text{traza}} = 1\times\frac{1414}{705} = \SI{2.01}{\kilo\metre},
 \qquad
 p_{\text{barrido}} = \frac{2{,}01}{\cos 65{,}5^\circ} = \SI{4.83}{\kilo\metre}.
\]
El píxel de <<1~km>> mide en el borde unos $2{,}0 \times 4{,}8$~km: casi diez veces más área que
en el nadir. El modelo plano se queda muy corto.
\end{ejemplo}

Este crecimiento tiene una consecuencia directa sobre la detección. Como se verá en el
capítulo~\ref{cap:deteccion}, un fuego se detecta cuando eleva lo suficiente la señal de
\emph{su} píxel. Si el píxel tiene diez veces más área, el mismo fuego queda diez veces más
diluido. \textbf{MODIS es mucho menos sensible en los bordes de su franja que en el centro.}

\section{El efecto pajarita}

MODIS no tiene un solo detector por banda sino una fila de diez, alineados en la dirección de
avance. En cada giro del espejo se registran diez líneas a la vez, que en el nadir cubren una
banda de 10~km de ancho en la dirección de avance; el satélite avanza justo esos 10~km entre un
giro y el siguiente, y las bandas encajan sin solaparse.

Pero hacia los bordes cada píxel crece también en la dirección de avance: en el extremo, las
diez líneas cubren unos 20~km en lugar de 10. El satélite sigue avanzando 10~km por giro, así
que cada barrido \textbf{pisa} la mitad del anterior y la mitad del siguiente. Visto en un mapa,
cada barrido tiene forma de corbata de lazo: estrecho en el centro y ancho en los extremos. De
ahí el nombre \termino{efecto pajarita} (\emph{bow-tie}), que muestra la
figura~\ref{fig:bowtie}.

\begin{figure}[ht]
\centering
\begin{tikzpicture}[x=0.85cm,y=1cm]
  \foreach \k/\c in {0/modis,1/fuego,2/bosque}{
    \fill[\c,opacity=0.28] (-6,\k*1.0-1.0)--(0,\k*1.0-0.5)--(6,\k*1.0-1.0)--(6,\k*1.0+1.0)--(0,\k*1.0+0.5)--(-6,\k*1.0+1.0)--cycle;
    \draw[\c,thick] (-6,\k*1.0-1.0)--(0,\k*1.0-0.5)--(6,\k*1.0-1.0)--(6,\k*1.0+1.0)--(0,\k*1.0+0.5)--(-6,\k*1.0+1.0)--cycle;
  }
  \node[font=\footnotesize,modis] at (-1.3,0) {barrido 1};
  \node[font=\footnotesize,fuego] at (0,1) {barrido 2};
  \node[font=\footnotesize,bosque] at (0,2) {barrido 3};
  \shade[inner color=yellow,outer color=fuego] (5.3,0.75) circle (0.13);
  \draw[<-,thick] (5.45,0.75)--(6.6,0.75) node[right,font=\footnotesize,align=left]{este fuego cae\\ en dos barridos};
  \draw[->,thick,ceniza] (-6.9,-0.8)--(-6.9,2.8) node[above,font=\footnotesize]{avance};
  \draw[<->] (-6,-1.5)--(6,-1.5) node[midway,below,font=\footnotesize]{franja de barrido (nadir en el centro)};
  \draw[<->] (0,-0.5)--(0,0.5);
  \node[font=\footnotesize,modis,fill=white,inner sep=1.5pt] at (1.3,-.25) {10 km};
  \draw[<->] (-6.25,-1)--(-6.25,1); \node[left,font=\footnotesize] at (-6.3,0) {20 km};
\end{tikzpicture}
\caption{El efecto pajarita en MODIS. Cada barrido registra una banda de 10~km de ancho en el
nadir que se ensancha hasta unos 20~km en los bordes. Barridos consecutivos se solapan lejos del
nadir, y un fuego situado allí queda registrado en dos de ellos.}
\label{fig:bowtie}
\end{figure}

\begin{atencion}[Focos duplicados dentro de una misma pasada]
Lejos del nadir, un mismo fuego puede aparecer como dos o más focos de MODIS con coordenadas
ligeramente distintas y exactamente la misma hora. No son dos fuegos. Si se cuentan <<puntos>>
sin más, los fuegos observados de lado pesan más que los observados desde arriba.
\end{atencion}

\section{La solución de VIIRS}

VIIRS se diseñó para corregir estos dos defectos. Usa dos ideas.

\subsection{Agregación de muestras}

Los detectores de VIIRS son, a lo largo del barrido, tres veces más finos de lo que el producto
final necesita. Cerca del nadir, el instrumento \emph{promedia tres muestras consecutivas} para
formar un píxel de 375~m. Más lejos, cuando cada muestra ya ha crecido, promedia solo dos. En la
parte exterior de la franja usa cada muestra tal cual. Hay así tres zonas de agregación:

\begin{center}
\begin{tabular}{lcc}
\toprule
Zona & Ángulo de barrido & Muestras promediadas \\
\midrule
Central & $0^\circ$ a $31{,}7^\circ$ & 3 \\
Intermedia & $31{,}7^\circ$ a $44{,}9^\circ$ & 2 \\
Exterior & $44{,}9^\circ$ a $56{,}3^\circ$ & 1 \\
\bottomrule
\end{tabular}
\end{center}

Cada vez que el píxel está a punto de crecer demasiado, se reduce el número de muestras y el
tamaño <<se reinicia>>. El resultado es la curva en dientes de sierra de la
figura~\ref{fig:tamano}: el píxel de 375~m de VIIRS no pasa de unos 800~m en el borde de la
franja, mientras que el de MODIS llega casi a 5~km.

\begin{figure}[ht]
\centering
\begin{tikzpicture}
\begin{axis}[width=13.5cm,height=7.4cm,xmin=0,xmax=1550,ymin=0,ymax=5.2,
  xlabel={Distancia al nadir sobre el suelo (km)},
  ylabel={Tamaño del píxel a lo largo del barrido (km)},
  ytick={0,1,2,3,4,5},
  legend style={at={(0.5,-0.2)},anchor=north,font=\small,draw=none,legend columns=2,/tikz/every even column/.append style={column sep=8pt}},
  grid=major,grid style={ceniza!15}]
\addplot[modis,very thick] table[x=dist,y=scan]{figuras/pixel_modis.dat};
\addplot[viirs,very thick] table[x=dist,y=scan]{figuras/pixel_viirs.dat};
\addplot[modis,thick,dashed] table[x=dist,y=track]{figuras/pixel_modis.dat};
\addplot[viirs,thick,dashed] table[x=dist,y=track]{figuras/pixel_viirs.dat};
\legend{MODIS a lo largo del barrido,VIIRS 375 m a lo largo del barrido,MODIS a lo largo de la traza,VIIRS 375 m a lo largo de la traza}
\draw[ceniza,dashed] (axis cs:523,0)--(axis cs:523,1.6);
\draw[ceniza,dashed] (axis cs:883,0)--(axis cs:883,1.6);
\node[font=\footnotesize,ceniza] at (axis cs:260,1.85) {3 muestras};
\node[font=\footnotesize,ceniza] at (axis cs:703,1.85) {2};
\node[font=\footnotesize,ceniza] at (axis cs:1200,1.85) {1 muestra};
\draw[->,ceniza] (axis cs:1300,4.3)--(axis cs:1175,4.75);
\node[font=\footnotesize,anchor=east] at (axis cs:1310,4.2) {borde de la franja de MODIS};
\end{axis}
\end{tikzpicture}
\caption{Tamaño del píxel según la distancia al nadir, calculado con las
ecuaciones~\eqref{eq:cenit} y~\eqref{eq:curvo}. El píxel de MODIS se multiplica casi por cinco
a lo largo del barrido; el de VIIRS, gracias a la agregación en tres zonas, se mantiene por
debajo de 0,8~km en toda la franja.}
\label{fig:tamano}
\end{figure}

\subsection{Eliminación del solape}

VIIRS también sufre el efecto pajarita, pero lo resuelve a bordo: en las zonas intermedia y
exterior, el propio instrumento descarta las líneas de cada barrido que quedarían duplicadas en
el barrido vecino. A esto se le llama \termino{eliminación de pajarita} (\emph{bow-tie
deletion}). Los datos de VIIRS llegan así casi sin repeticiones entre barridos de una misma
pasada.

\section{Las columnas \texttt{scan} y \texttt{track}}

Los datos de focos incluyen para cada detección el tamaño real del píxel en el que se hizo. La
columna \texttt{scan} es su dimensión a lo largo del barrido y \texttt{track} a lo largo de la
traza, ambas en kilómetros. En MODIS, \texttt{scan} va de 1,0 a 4,8 y \texttt{track} de 1,0 a
2,0. En VIIRS ambos se mantienen entre 0,3 y 0,8 aproximadamente.

Estas dos columnas son la manera más directa de saber \emph{desde dónde} se observó cada foco:
un valor de \texttt{scan} cercano a 1 en MODIS indica una observación casi vertical; un valor de
4 indica el borde de la franja.

\begin{ejemplo}{interpretar \texttt{scan} y \texttt{track}}
Un foco de MODIS tiene \texttt{scan} = 2,9 y \texttt{track} = 1,6. ¿Qué área cubre su píxel y a
qué distancia aproximada del nadir se observó? Un fuego de 1000~m$^2$, ¿qué fracción del píxel
ocupa aquí y qué fracción ocuparía en el nadir?
\solucion El área es aproximadamente $2{,}9 \times 1{,}6 = \SI{4.6}{\kilo\metre\squared}$. En
la figura~\ref{fig:tamano}, un tamaño de 2,9~km a lo largo del barrido corresponde a unos
860--880~km del nadir. El fuego ocupa $1000/\num{4.6e6} = \num{2.2e-4}$ del píxel, frente a
$1000/\num{1.0e6} = \num{1.0e-3}$ en el nadir: está casi cinco veces más diluido.
\end{ejemplo}

\begin{parareto}
Dos usos prácticos de estas columnas. Primero, como \textbf{filtro de calidad}: quedarse solo
con las detecciones de MODIS con \texttt{scan} menor que cierto valor (por ejemplo 2) elimina
la parte de la franja donde MODIS es poco sensible y donde aparecen duplicados. Segundo, como
\textbf{diagnóstico}: si la razón entre los conteos de VIIRS y de MODIS es mucho mayor para
detecciones de MODIS con \texttt{scan} alto, estás viendo directamente el efecto de este
capítulo. Cualquier filtro reduce el número de datos, así que hay que aplicarlo igual en todo el
periodo y documentarlo.
\end{parareto}

\section{Resolución espacial y tamaño mínimo detectable}

Queda una idea por fijar. Se suele decir que MODIS tiene <<1~km de resolución>> y VIIRS <<375
m>>. Eso describe el tamaño del píxel, no el tamaño del fuego más pequeño que pueden detectar.
Ambos detectan fuegos \emph{mucho} menores que su píxel, porque el fuego es tan brillante a
4~\um{} que una fracción diminuta del píxel basta para alterar su señal. Bajo buenas
condiciones, MODIS detecta de forma habitual fuegos de unos 1000~m$^2$ (una milésima de su
píxel) y, en condiciones muy favorables, de unos 100~m$^2$. VIIRS, con píxeles siete veces
menores en área, detecta fuegos proporcionalmente más pequeños. El capítulo~\ref{cap:deteccion}
explica el porqué con números.

\resumen

\begin{itemize}
  \item MODIS y VIIRS forman la imagen por \textbf{barrido transversal}: un espejo recorre
  líneas perpendiculares al avance.
  \item El \textbf{IFOV} es un ángulo fijo; la huella en el suelo crece con el ángulo de
  barrido. Con la Tierra curva el crecimiento es mucho mayor que con el modelo plano.
  \item El píxel de MODIS pasa de $1\times1$~km en el nadir a unos $2{,}0\times4{,}8$~km en el
  borde: diez veces más área y mucha menos sensibilidad.
  \item El \textbf{efecto pajarita} hace que barridos consecutivos se solapen lejos del nadir y
  produce focos duplicados en MODIS.
  \item VIIRS limita el crecimiento con \textbf{agregación en tres zonas} (su píxel de 375~m no
  supera 0,8~km) y elimina a bordo las líneas solapadas.
  \item Las columnas \texttt{scan} y \texttt{track} dan el tamaño real del píxel de cada foco.
\end{itemize}

\ejercicios

\begin{enumerate}
  \item Calcula el IFOV, en radianes y en grados, de la banda de 375~m de VIIRS a 824~km de
  altura. ¿Es mayor o menor que el de MODIS?
  \item Con el modelo plano, ¿a qué ángulo de barrido se duplica el tamaño del píxel a lo largo
  del barrido? ¿Y a lo largo de la traza?
  \item Repite el ejemplo del borde de la franja para $\theta = 40^\circ$ con el modelo curvo.
  ¿Cuántas veces mayor que en el nadir es el área del píxel?
  \item En la figura~\ref{fig:tamano}, estima el cociente entre el área del píxel de MODIS y el
  de VIIRS en el nadir y a 1100~km del nadir. ¿Dónde es mayor la ventaja de VIIRS?
  \item Explica por qué el efecto pajarita produce duplicados en los bordes de la franja pero no
  en el centro.
  \item Tienes un archivo de focos de MODIS de una región. Describe un procedimiento, usando
  solo las columnas de fecha, hora, coordenadas, \texttt{scan} y \texttt{track}, para identificar
  parejas de focos que probablemente sean el mismo fuego visto en dos barridos.
  \item Un colega dice: <<MODIS no puede ver fuegos menores de un kilómetro>>. Corrígelo en dos
  frases.
\end{enumerate}
